\documentclass[10pt,a4paper]{amsart}
\usepackage[margin=19mm]{geometry}
\usepackage{amsmath,amssymb,mathtools}
\usepackage[scr=rsfs,cal=euler]{mathalfa}
\usepackage{xfrac}
\usepackage[unicode,hidelinks,bookmarks=false]{hyperref}
\newcommand{\ie}{i.e.\ }
\newcommand{\cf}{cf.\ }
\newcommand{\resp}{respectively\ }
\newcommand{\Subsection}{\subsection}
\providecommand{\nobreakdash}{\nobreak\hbox{-}}
\setlength{\emergencystretch}{2em}
\multlinegap0pt
\usepackage[shortlabels]{enumitem}
\newtheorem{thm}{Theorem}
\newtheorem{prop}{Proposition}
\newcommand{\dx}[1][x]{%
	\ensuremath{
		{\mathrm{d}}{#1}
	}
}
\newcommand{\hr}{\le_{\mathrm{hr}}}
\newcommand{\st}{\le_{\mathrm{st}}}
\DeclareMathOperator{\supp}{supp}
\title{Stochastic Ordering under Weaker Likelihood-Ratio Shape Conditions}
\author{Z. Derbazi}
\date{Source: arXiv:2605.11973v1 (CC BY 4.0)}
\begin{document}
\maketitle

\noindent\emph{Source and licence.} Stochastic Ordering under Weaker Likelihood-Ratio Shape Conditions. Version arXiv:2605.11973v1 (CC BY 4.0), distributed by arXiv under the Creative Commons Attribution 4.0 International licence, http://creativecommons.org/licenses/by/4.0/, as recorded in the arXiv licence metadata for this exact version and shown on the abstract page https://arxiv.org/abs/2605.11973v1. Full licence notice: \texttt{licenses/arxiv-2605.11973-CC-BY-4.0.txt}.\par\smallskip

\noindent\textit{Editorial scope.} Counted unit: the article's prop prop:ylp-one-sided-score (source0.tex lines 368--375). The article's complete original proof of it is delivered with it (source0.tex lines 377--381, one \texttt{proof} environment of 5 lines). No mathematics is written, repaired or re-derived: the statement and the proof are the article's own lines, in the article's own notation, and no retained line is substituted or re-flowed.  Retained uncounted as the source-local context the proof needs: lines 92--95; lines 151--159.  Nothing was omitted from any retained span, so the package carries no omission record.  No retained line contains a citation, so the wrapper carries no bibliography and no bibliography entry.  No retained line contains a figure environment, an image-inclusion command or drawing code, so no image is delivered and the package declares no asset dependency.  Editorial formatting only, in addition: an AMS \texttt{amsart} preamble that restates the article's own declarations and macros used by the retained lines, namely the environments \texttt{thm}, \texttt{prop} on separate counters, together with the standard packages those lines need.  The retained environments print in this document as Theorem 1, Proposition 1 in order of appearance, whereas the article prints the same items with the numbering of its own full document.  The complete native source of the article as distributed in the arXiv e-print for this exact version is bundled unchanged as \texttt{sources/200394/source0.tex}.
\begin{equation}
    \label{def.ell}
    \ell(x):=\frac{\dx[]P/\dx[]\rho}{\dx[]Q/\dx[]\rho}, \qquad x\in S,
\end{equation}

\begin{thm}
\label{thm:unimodal-criterion}
Let \(P,Q\) be probability measures on \(E\), and let \(\ell\) be the likelihood ratio on $S:=\supp(P)\cup\supp(Q)$ as in~\eqref{def.ell}. If \(S\) is an interval with finite left endpoint \(x_*:=\inf S\), and \(\ell\) is unimodal on \(S\), then the following assertions are equivalent:
\begin{enumerate}[label=\textup{(\roman*)}]
\item $P\st Q$.
\item $P\hr Q$.
\item $\ell(x_*+)\ge 1$, where $\ell(x_*+)$ denotes the right limit of \(\ell\) at $x_*$.
\end{enumerate}
\end{thm}

\begin{prop}
\label{prop:ylp-one-sided-score}
Let $P,Q$ be probability laws on $E$ whose common support interval $J\subseteq E$ has finite left endpoint $x_0:=\min J$. If $\ell(x_0)\ge 1$ and the superlevel set
\[
A:=\{x\in J:\ell(x)\ge 1\}
\]
is an interval, then $P\st Q$. If, in addition, $\ell$ is nonincreasing on $J\setminus A$, then $P\hr Q$.
\end{prop}

\begin{proof}
Since $\ell(x_0)\ge 1$, the left endpoint belongs to $A$, so the interval assumption gives $A=[x_0,r]\cap J$ for some $r\in J\cup\{\sup J\}$. Hence $\ell\ge 1$ on $A$ and $\ell\le 1$ on $J\setminus A$, which means $f_P-f_Q$ has at most one sign change on $J$, from $+$ to $-$. The argument from~\textup{(iii)}\,$\Rightarrow$\,\textup{(i)} of Theorem~\ref{thm:unimodal-criterion} (using only the single sign change of $f_P-f_Q$, not the unimodality of $\ell$) gives $P\st Q$.

For the $\hr$ direction, set $T(x):=\bar F_P(x)/\bar F_Q(x)$ on $\{x\in J:\bar F_Q(x)>0\}$. On $A$ we have $\ell\ge 1$, so $f_P\ge f_Q$, and combining with $\bar F_P\le \bar F_Q$, established in the first part, gives $h_P\ge h_Q$. On $J\setminus A$, $\ell$ is nonincreasing and $\ell\le 1$. The bound $T(x)\le \ell(x)$ on $J\setminus A$ follows from the same tail-comparison proof of \textup{(i)}\,$\Rightarrow$\,\textup{(ii)} in Theorem~\ref{thm:unimodal-criterion}\textup{(i)}\,$\Rightarrow$\,\textup{(ii)}: for $x\in J\setminus A$ and $y\ge x$, monotonicity of $\ell$ on $J\setminus A$ gives $\ell(y)\le\ell(x)$, so $\bar F_P(x)\le \ell(x)\bar F_Q(x)$. That is, $T(x)\le \ell(x)$. The discrete and continuous monotonicity arguments of that proof then yield $T$ nonincreasing on $J\setminus A$. Hence, $h_P\ge h_Q$ holds also on $J\setminus A$.
\end{proof}

\end{document}
